% =====================================================================
% פרק 1 — מסגרת המחקר
% Owner: Agent 1 (INTRO_HISTORY_AGENT). Numbers aligned to comms/data_summary.json
% =====================================================================
\chapter{מסגרת המחקר}
\label{ch:research-framework}

\begin{managerbox}
מבצעי המימון הקבלניים הפכו בשנים \en{2023–2025} לכלי השיווק המרכזי בשוק הדירות החדשות, והם דוחים את רגע ההכרעה הפיננסית של אלפי משקי בית אל מועד המסירה. פרק זה מגדיר את הבעיה כמותית, מציג את שאלת המחקר ואת שלוש ההשערות, ומפרט כיצד כל השערה תיבחן. זוהי המסגרת להערכת החשיפה של בנק, יזם או רגולטור ל״קיר המסירות״ של \en{2026–2027}.
\end{managerbox}

\section{רקע: ממחזור ריבית למבצעי מימון}
\label{sec:ch1-background}

מחזור העלאות הריבית של בנק ישראל בשנים \en{2022–2023} (סעיף~\ref{sec:ch3-rates}) ייקר בבת אחת את המימון של רוכשי הדירות ושל היזמים, והיקף העסקאות ירד (סעיף~\ref{sec:ch3-demand}). במקום להוריד מחירים בגלוי, אימצו יזמים רבים \term{מבצעי מימון קבלניים}{\mbox{Developer Financing Promotions}}, ובראשם מבצעי \deal{20/80} ו־\deal{20/50/30} ו\term{הלוואות בלון}{Balloon/Bullet Loans} שהריבית עליהן ממומנת בידי היזם. באוגוסט \en{2026} דווחה הטבת מימון ב־\pct{30} ממכירות הקבלנים בשוק החופשי בחמשת אזורי המעקב של משרד האוצר \cite{mof_review_2026_08,nadlancenter_2026_08_promos}.

מבנים אלה דומים למנגנונים שקדמו למשבר הסאב־פריים בארצות הברית ולמשבר הנדל״ן בספרד: תשלום התחלתי נמוך באופן מלאכותי, מינוף יזמי גבוה ומלאי גדול של דירות לא־מכורות \cite{fcic2011,akin2014}. מכאן הצורך לבחון אם ישראל משחזרת את מנגנוני הכשל של \en{2008}, או שמא הבלמים המוסדיים שלה מסיטים את הסיכון למוקד אחר. איור~\ref{fig:research-funnel} מציג את הצמצום מן הזירה העולמית אל שאלת המחקר.

\begin{figure}[!htb]
\centering
\begin{english}
\begin{tikzpicture}[font=\small, lvl/.style={draw=white, line width=1.2pt},
  txt/.style={align=center, text=white, inner sep=0pt}]
% trapezoid levels: half-widths 6.2 -> 5.4 -> 4.6 -> 3.8 -> 3.0; explicit line breaks only
\path[lvl, fill=brandNavy]  (-6.2,0)    -- (6.2,0)    -- (5.4,-1.1) -- (-5.4,-1.1) -- cycle;
\path[lvl, fill=brandSteel] (-5.4,-1.1) -- (5.4,-1.1) -- (4.6,-2.2) -- (-4.6,-2.2) -- cycle;
\path[lvl, fill=brandNavy!85!brandRed] (-4.6,-2.2) -- (4.6,-2.2) -- (3.8,-3.3) -- (-3.8,-3.3) -- cycle;
\path[lvl, fill=brandRed]   (-3.8,-3.3) -- (3.8,-3.3) -- (3.0,-4.4) -- (-3.0,-4.4) -- cycle;
\node[txt] at (0,-0.55) {\he{\textbf{מחזורי אשראי ולקחי משבר הסאב־פריים}}\\\he{ארה״ב וספרד (פרקים \ref{ch:subprime}, \ref{ch:comparative})}};
\node[txt] at (0,-1.65) {\he{\textbf{שוק הדיור בישראל, \en{2022–2026}}}\\\he{ריבית, מחירים ומלאי (פרק \ref{ch:israel})}};
\node[txt] at (0,-2.75) {\he{\textbf{מבצעי קבלנים: \en{20/80}, \en{20/50/30}}}\\\he{אשראי צל (\en{H1}; פרק \ref{ch:methodology})}};
\node[txt] at (0,-3.85) {\he{\textbf{״קיר המסירות״ \en{2026–2027}}}\\\he{חוסן בנקים ויזמים (\en{H2}, \en{H3})}};
\end{tikzpicture}
\end{english}
\caption{משפך המחקר: ממחזורי האשראי העולמיים אל ״קיר המסירות״ הישראלי}
\label{fig:research-funnel}
\end{figure}

\section{הצגת הבעיה במונחים כמותיים}
\label{sec:ch1-problem}

טבלה~\ref{tab:ch1-problem} מסכמת את הבעיה: מצב נוכחי, מצב רצוי, הפער ועלותו. ערכים שמקורם בסימולציה מסומנים כתרחיש להמחשה.

\begin{table}[!htb]
\centering
\caption{הצגת הבעיה: מצב נוכחי, מצב רצוי, פער ועלות}
\label{tab:ch1-problem}
\footnotesize
\begin{tabularx}{\tabmaxwidth}{>{\raggedleft\arraybackslash}p{2.0cm} >{\hsize=.9\hsize}R >{\hsize=1.1\hsize}R}
\toprule
\hd{ממד} & \hd{מצב נוכחי (מאומת)} & \hd{מצב רצוי / פער ועלות} \\
\midrule
מלאי לא־מכור & \en{85,080} דירות חדשות לא־מכורות (ספטמבר \en{2026}); שיא של \en{86,290} בינואר \en{2026} \cite{cbs_newdwellings_2026q3,cbs_newdwellings} & מלאי התואם את קצב הספיגה. הפער: הון יזמי קשור הנושא עלויות מימון שוטפות \\
\midrule
שקיפות מחירים & מדד מחירי הדירות ירד ב־\pct{1.8} בשנה (יולי–אוגוסט \en{2026}) \cite{toi_housing_2026_10} & מחיר מדווח השווה למחיר הכלכלי. הפער: הנחה כלכלית גלומה של \en{10.9\%–12.8\%} לעסקת \deal{20/80} (תרחיש להמחשה, פרק~\ref{ch:methodology}) \\
\midrule
יכולת ההחזר של הרוכשים & הטבת מימון ב־\pct{30} ממכירות הקבלנים בשוק החופשי (חמישה אזורים, אוגוסט \en{2026}) \cite{mof_review_2026_08,nadlancenter_2026_08_promos} & יחס החזר במסירה מתחת ל־\pct{40}. הפער: בתרחיש בסיס \pct{36} מהרוכשים מעל הסף (סימולציה, פרק~\ref{ch:delivery-wall}) \\
\midrule
חשיפת המערכת הבנקאית & כ־\pct{39} מהאשראי העסקי מופנה לבינוי ולנדל״ן (\en{2025}) \cite{boi2026_recredit} & חשיפה מגודרת. העלות: הפסדי אשראי וצורך בהון נוסף אם הסיכון יעבור מהיזמים לבנקים \\
\bottomrule
\end{tabularx}
\end{table}

\noindent במשפט אחד: מבצעי המימון מחליפים כיום ירידת מחיר גלויה בהנחה סמויה, ודוחים לשנים \en{2026–2027} את מבחן יכולת ההחזר של רוכשים שמספרם הכולל לא פורסם (סעיף~\ref{sec:ch3-tam}). טרם ידוע אם הסיכון ייספג במאזני היזמים או יעבור אל המערכת הבנקאית.

\section{המסגרת התאורטית}
\label{sec:ch1-theory}

\paragraph{השערת חוסר היציבות של מינסקי.}\label{subsec:ch1-minsky}
\sloppy תקופות ממושכות של יציבות מולידות את חוסר היציבות הבא, לפי \term{השערת חוסר היציבות הפיננסית}{Financial Instability Hypothesis} \cite{minsky1986,minsky1992}. ב\term{מימון גידור}{Hedge Finance} התזרים מכסה קרן וריבית; ב\term{מימון ספקולטיבי}{Speculative Finance} הוא מכסה רק את הריבית; וב\term{מימון פונזי}{Ponzi Finance} הלווה תלוי בעליית ערך הנכס. ככל שמשקל שני האחרונים גדל, המערכת שבירה יותר, עד ל״רגע מינסקי״ \enb{Minsky Moment}. רוכש במבצע \deal{20/80} המסתמך על משכנתא עתידית בתנאים לא ידועים נמצא במבנה ספקולטיבי ואף פונזי, וכך גם יזם הממנן את ריבית הבלון מאשראי ותלוי במכירת המלאי.\par\fussy

\paragraph{מידע אסימטרי וסיכון מוסרי.}\label{subsec:ch1-asymmetric}\label{subsec:ch1-moral-hazard}
כאשר צד אחד אינו יכול להבחין באיכות הנכס או הלווה, המחיר משקף איכות ממוצעת, והמלווים נוטים לקצוב אשראי — \term{ברירה שלילית}{Adverse Selection} \cite{akerlof1970,stiglitz1981}. \term{סיכון מוסרי}{Moral Hazard} נוצר כאשר צד המוגן מהשלכות התנהגותו נוטל סיכון מוגבר \cite{holmstrom1979}. במבצעי המימון פועלים שני המנגנונים: שווי ההטבה אינו נכלל במחיר המדווח, ומבחן המשכנתא נדחה למסירה. הקצאת הסיכון בין הצדדים היא לב \en{H3}.

\paragraph{בנקאות צל ואשראי צל.}\label{subsec:ch1-shadow}
\term{בנקאות צל}{Shadow Banking} היא תיווך אשראי מחוץ למערכת המפוקחת או בדרכים העוקפות את מגבלותיה \cite{pozsar2010,gorton2012}. \term{אשראי צל}{Shadow Credit} מוגדר בעבודה זו כהסדר שמגדיל בפועל את מינוף משק הבית או היזם, אך אינו נרשם ככזה במסגרת הכללים. מבצעי \deal{20/80} עונים להגדרה: הרוכש מתחייב לרכישה מלאה, ומבחני \en{LTV/PTI} נדחים.


\section{שאלת המחקר והשערותיו}
\label{sec:ch1-rq}

\begin{researchquestion}
\ResearchQuestion
\end{researchquestion}

\noindent השאלה מתפרקת לשלוש שאלות משנה: \textbf{מנגנון} — האם המבצעים פועלים כאשראי צל המסווה ירידת מחירים ריאלית (\en{H1}); \textbf{חשיפת משקי הבית} — מה יעלה בגורלם של רוכשים בתשלומים נדחים בעת המסירה (\en{H2}); ו\textbf{מיקום הסיכון המערכתי} — במאזני הבנקים, כבארה״ב ובספרד, או במאזני היזמים ובמלאי (\en{H3}).

\label{sec:ch1-hypotheses}ניסוח ההשערות שלהלן מחייב, והוא חוזר מילה במילה בפרק~\ref{ch:delivery-wall}, שבו הן נבחנות.

\begin{researchquestion}[השערות המחקר]
\textbf{\en{H1} — אשראי צל וקשיחות מחירים מלאכותית:} מבצעי המימון הקבלניים (\deal{20/80}, \deal{20/50/30}) והלוואות הבלון תפקדו כאשראי צל \enb{Shadow Credit} שעקף את מגבלות שיעור המימון (\en{LTV}) ויחס ההחזר מההכנסה (\en{PTI}) של בנק ישראל, ויצרו קשיחות מלאכותית של מחירי המחירון, תוך הסוואת ירידה ריאלית במחיר הכלכלי של הדירה במונחי ערך נוכחי נקי (\en{NPV}).

\medskip
\textbf{\en{H2} — לחץ החזר בעת ״קיר המסירות״:} רמת ריבית המשכנתאות הגבוהה בתקופת ״קיר המסירות״ (\en{2026–2027}) תדחוף שיעור מהותי מרוכשי הדירות בתשלומים נדחים אל מעבר לסף יחס ההחזר של \pct{40} מההכנסה (\en{PTI}), שמעליו חל משקל סיכון מוגבר, ותוביל למצוקת מימון מחדש \enb{Refinancing Distress} ולביטולי עסקאות פוטנציאליים.

\medskip
\textbf{\en{H3} — הסטת הסיכון המערכתי ליזמים:} בשל הבלמים המוסדיים — חוק המכר (דירות), מנגנון הליווי הבנקאי הסגור, מגבלות \en{LTV/PTI} נוקשות והלוואות בחבות אישית מלאה \enb{Full Recourse} — הסיכון המערכתי בישראל מתרכז במאזני היזמים ובמלאי הדירות הלא־מכורות (למעלה מ־\en{85,000} יחידות) ולא בחדלות פירעון של הבנקים, כפי שאירע בארה״ב ובספרד ב־\en{2008}.
\end{researchquestion}

\section{היקף, הנחות יסוד ודרישות}
\label{sec:ch1-scope}

\paragraph{בתוך ההיקף.} שוק הדירות החדשות בישראל בשנים \en{2022–2026}; מבצעי \deal{20/80}, \deal{20/50/30} והלוואות בלון במימון היזם; רוכשים פרטיים; יזמים בליווי בנקאי סגור; חמש הקבוצות הבנקאיות הגדולות. מקרי ההשוואה: ארה״ב (\en{2004–2009}) וספרד (\en{2004–2013}). \textbf{מחוץ להיקף:} שכירות, דיור ציבורי ו״מחיר למשתכן״, נדל״ן מסחרי והשפעות מאקרו מעבר לערוץ הריבית.

\paragraph{הנחות יסוד.} (\en{1}) נתוני בנק ישראל, הלמ״ס ומשרד האוצר מהימנים, בכפוף לכך שספירת המלאי של הלמ״ס כוללת גם דירות שבנייתן טרם החלה; (\en{2}) מחיר העסקה המדווח הוא מחיר המחירון, ללא שווי ההטבה המימונית; (\en{3}) פרמטרים שאינם נצפים מכוילים ומוצגים כתרחיש (פרק~\ref{ch:methodology}); (\en{4}) תנאי הרגולציה הם אלה שבתוקף בנובמבר \en{2026}, ובהם הוראת השעה שתוקפה עד \en{31.12.2026} \cite{boi_circular2816}.

\paragraph{דרישות ושיטת אימות.}\label{sec:ch1-requirements} כל השערה מתורגמת לדרישה מדידה עם ערך יעד ושיטת אימות (טבלה~\ref{tab:ch1-requirements}). ספי ההכרעה נקבעו לפני הרצת המודלים, כדי למנוע התאמה בדיעבד.

\begin{table}[!htb]
\centering
\caption{דרישות המחקר: מזהה, דרישה, ערך יעד ושיטת אימות}
\label{tab:ch1-requirements}
\footnotesize
\begin{tabularx}{\tabmaxwidth}{>{\raggedleft\arraybackslash}p{1.95cm} >{\hsize=.9\hsize}R >{\hsize=1.15\hsize}R >{\hsize=.95\hsize}R}
\toprule
\hd{מזהה} & \hd{דרישה} & \hd{ערך יעד / סף הכרעה} & \hd{שיטת אימות} \\
\midrule
\mbox{\en{REQ-H1-01}} & אומדן ההנחה הכלכלית הגלומה בעסקאות \deal{20/80} ו־\deal{20/50/30} & \en{H1} נתמכת אם ההנחה עולה על \pct{5} בכל שיעורי ההיוון \en{5\%–6\%} & מודל \en{NPV}, פרק~\ref{ch:methodology} \\
\midrule
\mbox{\en{REQ-H1-02}} & השוואת מדד המחירים הרשמי למדד מתוקן \en{NPV} & סדרה חודשית \en{2021–2026}; פער חיובי מתמשך & מדד מתוקן מול נתוני הלמ״ס \\
\midrule
\mbox{\en{REQ-H2-01}} & התפלגות יחס ההחזר במסירה בשלושה תרחישי ריבית & \en{H2} נתמכת אם בתרחיש הבסיס \pct{20} או יותר מהרוכשים מעל \pct{40} & סימולציית מונטה קרלו, \mbox{\en{N = 10,000}}, פרק~\ref{ch:delivery-wall} \\
\midrule
\mbox{\en{REQ-H3-01}} & מבחן קיצון לפרויקט מייצג: ירידת מחירים \en{0\%–20\%} ועיכוב \en{0–24} חודשים & \en{H3} נתמכת אם בתרחיש החמור יחס הכיסוי של היזם יורד אל מתחת ל־\en{1} \enb{DSCR < 1}, בעוד ההפסד הבנקאי המוערך נמוך מ־\pct{10} מההון העצמי של המערכת & רשת \en{81} תאים ואומדן הפסד, פרק~\ref{ch:delivery-wall} \\
\midrule
\mbox{\en{REQ-CMP-01}} & השוואה רב־ממדית ארה״ב / ספרד / ישראל & \en{8} פרמטרים לפחות, כל אחד עם מקור & טבלת השוואה, פרק~\ref{ch:comparative} \\
\midrule
\mbox{\en{REQ-DAT-01}} & אמינות הנתונים & \pct{100} מהערכים העובדתיים עם מקור מאומת & הצלבה מול \en{data\_summary.json} \\
\bottomrule
\end{tabularx}
\end{table}

\paragraph{מבנה הדוח.}\label{sec:ch1-structure} מבנה הדוח עוקב אחר המשפך שבאיור~\ref{fig:research-funnel}; פרק~\ref{ch:delivery-wall} מכריע בהשערות, ופרק~\ref{ch:risks-gantt} מציג את תוכנית העבודה והסיכום. הפרק הבא פותח במשבר \en{2008} עצמו ובמה שהספרות למדה עליו.
